% =============================================================================
% 3.3 Performances
% =============================================================================
\section{Performances et évaluation}
\label{sec:ch3-perfs}

Les mesures ci-dessous portent sur un corpus de dossiers d'appels d'offres
et de pièces patrimoniales fournis par SITINFRA Cameroun, complété par des
documents publics de marchés. Elles constituent une première campagne, à
interpréter comme une démonstration de faisabilité plutôt que comme une
évaluation définitive.

\subsection{Performance de l'ingestion}

\begin{btptable}{Temps d'ingestion observés}
\label{tab:perf-ingestion}
\begin{tabularx}{\linewidth}{@{}C{2.4cm} L{3.6cm} Y@{}}
\toprule
\rowcolor{btppale}
\btpth{Taille} & \btpth{Type} & \btpth{Durée observée} \\
\midrule
20 pages  & PDF numérique & 8 à 12\,s \\
50 pages  & PDF numérique & 15 à 25\,s \\
20 pages  & Scan          & 30 à 45\,s \\
50 pages  & Scan          & 60 à 90\,s \\
\bottomrule
\end{tabularx}
\end{btptable}

La cible méthodologique (moins de soixante secondes pour cinquante pages
numériques) est tenue. Elle ne l'est pas pour les scans de même volume,
ce qui confirme que la reconnaissance optique reste le goulot
d'étranglement de l'horloge A. Ce dépassement n'affecte pas l'horloge B,
précisément parce que les deux horloges sont découplées.

\subsection{Performance de la recherche}

\begin{btptable}{Comparaison des approches de recherche}
\label{tab:perf-recherche}
\begin{tabularx}{\linewidth}{@{}L{4.0cm} C{2.4cm} C{2.6cm} C{2.2cm}@{}}
\toprule
\rowcolor{btppale}
\btpth{Approche} & \btpth{Rappel@20} & \btpth{Précision@5} & \btpth{Temps} \\
\midrule
Dense seul            & $0{,}72$ & $0{,}45$ & $0{,}3$\,s \\
BM25 seul             & $0{,}68$ & $0{,}38$ & $0{,}1$\,s \\
Hybride (RRF)         & $0{,}78$ & $0{,}52$ & $0{,}4$\,s \\
Hybride et ré-ordonnancement & $0{,}81$ & $0{,}61$ & $0{,}6$\,s \\
\bottomrule
\end{tabularx}
\end{btptable}

L'hypothèse H1 est soutenue~: le hybride gagne $0{,}06$ de rappel à 20
sur le dense seul, soit un peu plus de huit pour cent relatifs, proche
de la cible de dix pour cent. L'hypothèse H3 l'est également~: le
ré-ordonnancement gagne $0{,}09$ de précision à 5 sur la seule fusion,
soit environ dix-sept pour cent relatifs, légèrement sous la cible de
vingt pour cent, mais dans le bon sens et à un surcoût de latence
acceptable ($0{,}2$\,s).

\subsection{Performance de la génération}

\begin{btptable}{Scores de génération (moyenne sur cinq sections, trois évaluateurs)}
\label{tab:perf-generation}
\begin{tabularx}{\linewidth}{@{}L{4.4cm} C{3.0cm}@{}}
\toprule
\rowcolor{btppale}
\btpth{Métrique} & \btpth{Score} \\
\midrule
Fidélité au contexte           & $0{,}87$ \\
Ancrage des affirmations       & $0{,}82$ \\
Exhaustivité                   & $0{,}78$ \\
Taux de validation sans retouche & $72$\,\% \\
\bottomrule
\end{tabularx}
\end{btptable}

Les cibles de fidélité ($> 0{,}85$) et de taux de validation ($> 70$\,\%)
sont atteintes. L'ancrage ($0{,}82$) dépasse le seuil d'alerte de $0{,}70$
mais laisse une marge de progrès, principalement sur les faits chiffrés
dont la citation est parfois omise par le modèle malgré la consigne.

\subsection{Coûts}

\begin{btptable}{Coûts estimés pour dix dossiers par mois}
\label{tab:couts}
\begin{tabularx}{\linewidth}{@{}L{5.2cm} Y@{}}
\toprule
\rowcolor{btppale}
\btpth{Poste} & \btpth{Ordre de grandeur} \\
\midrule
Plongements (BGE-M3, local)     & Négligeable une fois le modèle chargé \\
Modèle de langage local (Ollama) & Nul en licence, coût matériel déjà amorti \\
Modèle de langage par API (option) & 10 à 30\,€ \\
Stockage objet et bases         & 5 à 10\,€ \\
\bottomrule
\end{tabularx}
\end{btptable}

L'exécution locale du modèle de langage, lorsqu'un accélérateur est
disponible, maintient le coût marginal proche de zéro et évite de faire
sortir des pièces confidentielles vers un prestataire externe. C'est un
argument de sécurité autant que de budget.
